% 2020 新高考Ⅱ卷 11：正弦型函数 y=sin(2x+2π/3) 的图象，x∈[−π/3, 7π/6]
% 绘制依据：源 PDF 原图。形态锚点：最大值在 x=−π/12（y 轴稍左）、x=π/6 处下穿 x 轴、
% 最小值 x=5π/12、x=2π/3 处上穿、x=11π/12 再现波峰；区间两端恰为零点（跨 1.5 个周期）。
% y 轴只标刻度 1（置于 (0,1) 左上方以避开紧邻 y 轴的波峰）；
% x 轴只标 π/6 与 2π/3，用手工节点分置左下 / 右下以避让曲线；O 标原点左下。
\begin{tikzpicture}
\begin{axis}[width=10.6cm, height=5.6cm, font=\small,
  axis lines=middle, xlabel={$x$}, ylabel={$y$},
  xmin=-1.35, xmax=4.05, ymin=-1.45, ymax=1.62,
  xtick=\empty, ytick=\empty,
  every axis x label/.style={at={(current axis.right of origin)}, anchor=west},
  every axis y label/.style={at={(current axis.above origin)}, anchor=south}]
  \addplot[domain=-1.0472:3.6652, samples=250, thick] {sin(deg(2*x+2*pi/3))};
  \draw (axis cs:-0.06,1) -- (axis cs:0.06,1);
  \node[above left, inner sep=1pt] at (axis cs:0,1) {$1$};
  \node[below left, inner sep=1pt] at (axis cs:0.5236,0) {$\frac{\pi}{6}$};
  \node[below right, inner sep=1pt] at (axis cs:2.0944,0) {$\frac{2\pi}{3}$};
  \node[below left, inner sep=1pt] at (axis cs:0,0) {$O$};
\end{axis}
\end{tikzpicture}
